\documentclass[a4paper, 12pt, one column, aas_macros]{article}

%% Language and font encodings. This says how to do hyphenation on end of lines.
\usepackage[english]{babel}
\usepackage[utf8x]{inputenc}
\usepackage[T1]{fontenc}
\usepackage{aas_macros}

%% Sets page size and margins. You can edit this to your liking
\usepackage[top=1.3cm, bottom=2.0cm, outer=2.5cm, inner=2.5cm, heightrounded,
marginparwidth=1.5cm, marginparsep=0.4cm, margin=2.5cm]{geometry}

%% Useful packages
\usepackage{graphicx} %allows you to use jpg or png images. PDF is still recommended
\usepackage[colorlinks=False]{hyperref} % add links inside PDF files
\usepackage{amsmath}  % Math fonts
\usepackage{amsfonts} %
\usepackage{amssymb}  %

%% Citation package
\usepackage[authoryear]{natbib}
\bibliographystyle{mnras}
\setcitestyle{authoryear,open={(},close={)}}

\title{Add title here}
\author{Author name}

\begin{document}
\maketitle

\begin{abstract}
Add abstract here
\end{abstract}

\section{Introduction}
\label{intro}

This is an example citation. See file refs.bib by clicking on ``Project'' on the top-left. 

You can cite using \citet{Claydon17} to make a reference in the text. Or \citep{Claydon17} for a reference in parenthesis. 

\section{Something}

You can also reference figures and sections using their label name.

For instance, Figure \ref{frog} shows a picture of a frog and in Section \ref{intro} we show the introduction.

\begin{figure}
\begin{center}
\includegraphics[width=0.2\textwidth]{frog.jpg}
\caption{Insert caption here.}
\label{frog}
\end{center}
\end{figure}

\subsection{Sub-something}

Add content

\bibliography{refs}
\end{document}